\documentclass[a4paper,12pt]{article}
\usepackage[T2A]{fontenc}
\usepackage[utf8]{inputenc}
\usepackage[russian]{babel}
\usepackage{amsmath,amssymb,mathtools}
\usepackage{helvet}
\renewcommand{\familydefault}{\sfdefault}
\usepackage{geometry}
\geometry{top=23mm,bottom=23mm,left=24mm,right=24mm}
\usepackage{microtype}
\usepackage{setspace}
\usepackage{parskip}
\usepackage{array,tabularx,booktabs}
\usepackage{enumitem}
\usepackage{needspace}
\usepackage{xcolor}
\usepackage{tikz}
\usetikzlibrary{arrows.meta,positioning,shapes.geometric,calc}
\usepackage{pgfplots}
\pgfplotsset{compat=1.18}
\usepackage{fancyhdr}
\usepackage{hyperref}
\definecolor{accent}{HTML}{008080}
\definecolor{darkaccent}{HTML}{005757}
\definecolor{textgray}{HTML}{343B3D}
\definecolor{muted}{HTML}{627476}
\hypersetup{colorlinks=true,linkcolor=darkaccent,pdfauthor={Лёвин Артём Александрович},pdftitle={Чек-лист. Пересечение графиков и принадлежность точки}}
\setstretch{1.21}
\setlength{\parindent}{0pt}
\setlength{\parskip}{.62em}
\setlist[enumerate]{leftmargin=*,itemsep=.5em,topsep=.45em}
\setlist[itemize]{leftmargin=*,itemsep=.23em,topsep=.38em}
\pagestyle{fancy}
\fancyhf{}
\fancyhead[L]{\small\color{muted}Пересечение графиков и точки}
\fancyhead[R]{\small\color{muted}Кирилл Зиновьев}
\fancyfoot[C]{\small\color{muted}\thepage}
\renewcommand{\headrulewidth}{.35pt}
\renewcommand{\headrule}{\hbox to\headwidth{\color{accent}\leaders\hrule height\headrulewidth\hfill}}
\setlength{\headheight}{25pt}
\newcommand{\secname}[1]{\needspace{6\baselineskip}\vspace{.8em}{\Large\color{darkaccent}#1}\par\vspace{-.4em}{\color{accent}\rule{\linewidth}{.45pt}}\vspace{.10em}}
\newcommand{\subname}[1]{\needspace{3\baselineskip}{\color{darkaccent}\bfseries #1}\par\vspace{-.35em}}
\newcommand{\idea}[1]{\noindent{\color{darkaccent}\bfseries Главная идея.}\enspace #1}
\newcommand{\warn}[1]{\noindent{\color{darkaccent}\bfseries Проверьте.}\enspace #1}
\newcommand{\oneexercise}[2]{\item \textbf{#1} #2}
\begin{document}
\begin{titlepage}
\thispagestyle{empty}
\begin{center}
\vspace*{25mm}
{\large\color{darkaccent} МАТЕМАТИКА}\par
\vspace{21mm}
{\LARGE\color{textgray} Чек-лист}\par
\vspace{14mm}
{\LARGE\color{accent} Пересечение графиков\\[2mm] и принадлежность точки прямой}\par
\vspace{19mm}
{\large\color{textgray} 7 класс}\par
\vfill
{\large\color{textgray} Кирилл Зиновьев}\par
\vspace{14mm}
{\large\color{textgray} 10 ноября 2026 г.}\par
\vspace{15mm}
{\small\color{textgray} Лёвин Артём Александрович\\[2mm]
эксклюзивно для Кирилла Зиновьева}\par
\vspace{18mm}
\begin{tikzpicture}[x=1cm,y=1cm]
  \draw[accent,line width=.9pt] (0,0) .. controls (1.3,.65) and (2.5,-.62) .. (4.0,-.08)
    .. controls (5.2,.54) and (6.6,-.52) .. (8,0);
\end{tikzpicture}
\end{center}
\end{titlepage}
\setcounter{page}{1}

\secname{1. Координаты точки: абсцисса и ордината}
Каждой точке на координатной плоскости соответствуют \textbf{две координаты}:
\[
\boxed{T(x;y)}
\]
\begin{itemize}
  \item \textbf{Абсцисса} --- число \(x\), координата по горизонтальной оси \(Ox\).
  \item \textbf{Ордината} --- число \(y\), координата по вертикальной оси \(Oy\).
\end{itemize}
Всегда записывайте сначала \(x\), затем \(y\). Например, в точке \(A(2;5)\) абсцисса равна \(2\), ордината равна \(5\).

\secname{2. Вспоминаем: прямая по двум точкам}
График линейной функции \(y=kx+b\) --- прямая. Чтобы построить её, достаточно выбрать \textbf{два различных значения} \(x\), вычислить соответствующие \(y\) и провести прямую через полученные точки.

\subname{Пример: \(y=2x+5\)}
Возьмём \(x=0\) и \(x=1\). Подставим в формулу:
\[
\begin{aligned}
x=0&:\quad y=2\cdot0+5=5,\\
x=1&:\quad y=2\cdot1+5=7.
\end{aligned}
\]
Получили точки \(A(0;5)\) и \(B(1;7)\). Запишем данные в таблицу:
\[
\begin{array}{c|cc}
 & A & B\\ \hline
x & 0 & 1\\
y & 5 & 7
\end{array}
\]
\idea{Значения \(x\) можно выбрать самостоятельно. Значения \(y\) обязательно вычисляются подстановкой в данную формулу.}

\secname{3. Что такое точка пересечения графиков}
Точка пересечения принадлежит \textbf{обоим графикам одновременно}. Поэтому при её абсциссе \(x\) обе формулы дают \textbf{одно и то же значение} \(y\).

Для функций \(y_1=2x+5\) и \(y_2=-x+7\) можно изобразить обе прямые на одной плоскости. Их общая точка обозначена \(T\).

\begin{center}
\begin{tikzpicture}
\begin{axis}[
 width=11.0cm,height=7.0cm,
 xmin=-.9,xmax=2.1,ymin=3.0,ymax=9.1,
 axis lines=middle,
 xlabel={$x$},ylabel={$y$},
 xtick={-0.5,0,0.5,1,1.5,2},
 ytick={3,4,5,6,7,8,9},
 grid=major,major grid style={draw=gray!25,line width=.2pt},
 tick label style={font=\small},
 axis line style={black!65},
 clip=false,legend style={at={(1.02,.98)},anchor=north west,draw=none,font=\small},
 width=10.5cm
]
\addplot[accent,thick,domain=-.7:2,samples=2] {2*x+5};
\addlegendentry{$y=2x+5$}
\addplot[orange!70!black,thick,domain=-.7:2,samples=2] {-x+7};
\addlegendentry{$y=-x+7$}
\addplot[only marks,mark=*,mark size=2.4pt,black] coordinates {(0.6666667,6.3333333)};
\node[anchor=south east,font=\small,fill=white,inner sep=1.7pt] at (axis cs:.6666667,6.3333333) {$T\bigl(\frac23;\frac{19}3\bigr)$};
\end{axis}
\end{tikzpicture}
\end{center}
График помогает увидеть место пересечения. Но координаты, считанные с рисунка, часто приблизительны. \textbf{Для точного ответа используйте уравнение.}

\secname{4. Точное нахождение точки пересечения}
\idea{В точке пересечения \(y_1=y_2\). Приравняйте правые части формул и найдите \(x\). Затем подставьте найденный \(x\) в любую исходную формулу и получите \(y\).}

\subname{Пример 1. \(y_1=2x+5\), \(y_2=-x+7\)}
\textbf{Шаг 1.} Приравниваем выражения:
\[
2x+5=-x+7.
\]
\textbf{Шаг 2.} Собираем неизвестные слева, числа справа:
\[
2x+x=7-5,\qquad 3x=2,\qquad \boxed{x=\frac23}.
\]
Это \textbf{абсцисса} точки пересечения.

\textbf{Шаг 3.} Подставляем \(x=\frac23\) во вторую формулу:
\[
y=-\frac23+7=-\frac23+\frac{21}{3}=\boxed{\frac{19}{3}}.
\]
Это \textbf{ордината} точки пересечения. Обратите внимание: результат \textbf{положительный}, так как из \(7\) вычитается всего \(\frac23\).

\textbf{Ответ:} \(\displaystyle \boxed{T\left(\frac23;\frac{19}{3}\right)}\).

\textbf{Проверка по первой формуле:}
\[
2\cdot\frac23+5=\frac43+\frac{15}{3}=\frac{19}{3}.
\]
Обе формулы дают одинаковое значение. Координаты найдены правильно.

\subname{Пример 2. \(y_1=x+1\), \(y_2=-x+2\)}
\[
\begin{aligned}
x+1&=-x+2,\\
2x&=1,\\
x&=\frac12=0{,}5,\\
y&=0{,}5+1=1{,}5.
\end{aligned}
\]
\textbf{Ответ:} \(\boxed{T(0{,}5;1{,}5)}\).

\warn{Если требуется только абсцисса, достаточно найти \(x\). Если требуются координаты точки, обязательно найдите и \(x\), и \(y\).}

\clearpage
\thispagestyle{fancy}
\vspace*{7mm}
\begin{center}
{\Large\color{darkaccent}Алгоритм: точка пересечения двух прямых}\par
\vspace{13mm}
\begin{tikzpicture}[
  node distance=12mm,
  every node/.style={font=\small,align=center},
  startstop/.style={ellipse,draw=darkaccent,thick,minimum width=42mm,minimum height=11mm,inner sep=3pt},
  process/.style={rounded corners=2mm,rectangle,draw=accent,thick,text width=95mm,minimum height=13mm,inner sep=4pt},
  arrow/.style={-{Latex[length=3mm]},thick,draw=darkaccent}
]
\node[startstop] (start) {Даны две прямые\\$y=f(x)$ и $y=g(x)$};
\node[process,below=of start] (equal) {Приравняйте правые части:\\$f(x)=g(x)$};
\node[process,below=of equal] (findx) {Решите уравнение;\\найдите абсциссу $x_0$};
\node[process,below=of findx] (findy) {Подставьте $x_0$ в любую формулу;\\найдите ординату $y_0$};
\node[process,below=of findy] (verify) {Проверьте: $f(x_0)=g(x_0)=y_0$};
\node[startstop,below=of verify] (end) {Запишите $T(x_0;y_0)$};
\draw[arrow] (start)--(equal);
\draw[arrow] (equal)--(findx);
\draw[arrow] (findx)--(findy);
\draw[arrow] (findy)--(verify);
\draw[arrow] (verify)--(end);
\end{tikzpicture}
\end{center}
\clearpage

\secname{5. Принадлежит ли точка графику}
Пусть задана функция \(y=f(x)\) и точка \(A(x_0;y_0)\). Чтобы узнать, лежит ли точка на графике, необходимо подставить \(x_0\) в формулу и сравнить результат с \(y_0\):
\[
\boxed{A(x_0;y_0)\text{ лежит на графике}\quad\Longleftrightarrow\quad f(x_0)=y_0.}
\]

\subname{Пример 3. Функция \(y=x+1\)}
\textbf{Точка \(J(100;101)\).} Вычислим значение функции при \(x=100\):
\[
f(100)=100+1=101.
\]
Расчётное значение совпало с ординатой точки: \(101=101\). \textbf{Точка \(J\) лежит на графике.}

\textbf{Точка \(Q(100;99)\).} При том же \(x=100\) формула по-прежнему даёт \(101\), а ордината точки равна \(99\). Поскольку \(101\ne99\), \textbf{точка \(Q\) на графике не лежит.}

\subname{Пример 4. Функция \(y=2x-3\)}
Проверим две точки, \(A(0;-3)\) и \(B(1;5)\):
\[
\begin{array}{@{}ll@{}}
A:&2\cdot0-3=-3\quad\text{--- совпало, точка лежит;}\\[4pt]
B:&2\cdot1-3=-1\ne5\quad\text{--- не совпало, не лежит.}
\end{array}
\]
\idea{Для проверки принадлежности \textbf{строить прямую не требуется}. Достаточно одной подстановки на каждую точку.}

\secname{6. Типичные ошибки и самопроверка}
\begin{itemize}
  \item \textbf{Путают координаты:} \(x\) --- абсцисса, \(y\) --- ордината; всегда записывайте \((x;y)\).
  \item \textbf{Путают способы:} при построении выбирают \(x\) и вычисляют \(y\); при нахождении пересечения приравнивают две формулы.
  \item \textbf{Теряют знак при переносе:} из \(2x+5=-x+7\) получается \(3x=2\), поскольку \(7-5=2\).
  \item \textbf{Забывают вторую координату:} после нахождения \(x\) вычислите \(y\), если нужны координаты точки.
  \item \textbf{Ошибаются с дробями:} \(7-\frac23=\frac{21}3-\frac23=\frac{19}3\), знак результата положительный.
  \item \textbf{Проверяют лишь абсциссу:} для точки \(A(x_0;y_0)\) нужно сравнить \(f(x_0)\) именно с \(y_0\).
\end{itemize}

\clearpage
\thispagestyle{fancy}
\vspace*{7mm}
\begin{center}
{\Large\color{darkaccent}Алгоритм: проверка принадлежности точки}\par
\vspace{13mm}
\begin{tikzpicture}[
  node distance=13mm,
  every node/.style={font=\small,align=center},
  startstop/.style={ellipse,draw=darkaccent,thick,minimum width=42mm,minimum height=11mm,inner sep=3pt},
  process/.style={rounded corners=2mm,rectangle,draw=accent,thick,text width=88mm,minimum height=14mm,inner sep=4pt},
  decision/.style={diamond,aspect=2.0,draw=darkaccent,thick,text width=42mm,inner sep=2pt},
  result/.style={rounded corners=2mm,rectangle,draw=accent,thick,text width=39mm,minimum height=13mm,inner sep=3pt},
  arrow/.style={-{Latex[length=3mm]},thick,draw=darkaccent}
]
\node[startstop] (start) {Даны $y=f(x)$\\и точка $A(x_0;y_0)$};
\node[process,below=of start] (sub) {Подставьте $x_0$ в формулу;\\вычислите $f(x_0)$};
\node[decision,below=14mm of sub] (check) {$f(x_0)=y_0$?};
\node[result,below left=17mm and 13mm of check] (yes) {Точка лежит\\на графике};
\node[result,below right=17mm and 13mm of check] (no) {Точка не лежит\\на графике};
\draw[arrow] (start)--(sub);
\draw[arrow] (sub)--(check);
\draw[arrow] (check.south west) -| node[pos=.24,above left] {да} (yes.north);
\draw[arrow] (check.south east) -| node[pos=.24,above right] {нет} (no.north);
\end{tikzpicture}
\end{center}
\clearpage

\secname{Тренировочный блок}
\textbf{Путь А: 10 заданий.} Выполняйте их последовательно, от простого к более сложному. В заданиях на пересечение используйте расчётный способ; для проверки принадлежности график строить не требуется.

\subname{Базовый уровень}
\begin{enumerate}[label=\textbf{\arabic*.},itemsep=.72em]
\item Для функции \(y=2x+1\) найдите \(y\) при \(x=0\) и \(x=2\). Запишите координаты двух полученных точек.
\item Для функции \(y=-3x+4\) найдите значения \(y\) при \(x=0\) и \(x=1\). Запишите координаты точек.
\item Проверьте, принадлежит ли точка \(A(0;4)\) графику \(y=3x+4\).
\item Проверьте, принадлежит ли точка \(B(2;1)\) графику \(y=-x+3\).

\subname{Средний уровень}
\item Для графика \(y=2x-3\) определите, какие точки принадлежат прямой: \(A(0;-3)\), \(B(1;-1)\), \(C(2;2)\).
\item Найдите координаты точки пересечения прямых \(y=x+2\) и \(y=-x+6\).
\item Найдите координаты точки пересечения прямых \(y=3x-1\) и \(y=x+5\).

\subname{Уровень выше среднего}
\item Найдите координаты точки пересечения прямых \(y=4x+1\) и \(y=-2x+4\). Ответ разрешается записать дробями.
\item Найдите координаты точки пересечения прямых \(y=3x-4\) и \(y=-2x+5\). Проверьте результат подстановкой в обе формулы.
\item Найдите координаты точки пересечения прямых \(y=\frac12x+2\) и \(y=-x+5\). Проверьте, принадлежит ли найденная точка обеим прямым.
\end{enumerate}

\clearpage
\secname{Ответы к тренировочному блоку}
\renewcommand{\arraystretch}{1.38}
\begin{center}
\begin{tabularx}{\linewidth}{@{}>{\bfseries}p{12mm}X@{}}
\toprule
№ & Ответ\\
\midrule
1 & \(y(0)=1\), \(y(2)=5\); точки \((0;1)\), \((2;5)\).\\
2 & \(y(0)=4\), \(y(1)=1\); точки \((0;4)\), \((1;1)\).\\
3 & Да: \(3\cdot0+4=4\).\\
4 & Да: \(-2+3=1\).\\
5 & Принадлежат точки \(A\) и \(B\); точка \(C\) не принадлежит.\\
6 & \(T(2;4)\).\\
7 & \(T(3;8)\).\\
8 & \(\displaystyle T\left(\frac12;3\right)\).\\
9 & \(\displaystyle T\left(\frac95;\frac75\right)\).\\
10 & \(T(2;3)\); принадлежит обеим прямым.\\
\bottomrule
\end{tabularx}
\end{center}
\vspace{1.7em}
\secname{Итоговая памятка}
\renewcommand{\arraystretch}{1.28}
\begin{center}
\begin{tabularx}{\linewidth}{@{}p{.37\linewidth}X@{}}
\toprule
Что требуется & Как действовать\\
\midrule
Построить прямую & Возьмите два разных \(x\), найдите два \(y\), отметьте точки и проведите прямую.\\
Найти абсциссу пересечения & Приравняйте формулы: \(f(x)=g(x)\), найдите \(x\).\\
Найти обе координаты & После \(x\) подставьте его в одну формулу и найдите \(y\).\\
Проверить точку \(A(x_0;y_0)\) & Вычислите \(f(x_0)\), сравните с \(y_0\).\\
Проверить ответ & Подставьте найденный \(x\) в обе формулы; значения \(y\) должны совпасть.\\
\bottomrule
\end{tabularx}
\end{center}
\end{document}
